\section{Related Work}

The challenge of strategic bidding in Real-Time Bidding (RTB) markets resides at the intersection of high-frequency sequential decision-making and adversarial multi-agent dynamics. While traditional approaches have evolved from simple heuristic controllers to sophisticated deep reinforcement learning (DRL) agents, they consistently struggle with the latent non-stationarity of competitor strategies and the inherent information asymmetry that leads to the Winner's Curse.

\subsection{Reinforcement Learning in Real-Time Bidding}
Early methodologies in RTB primarily utilized linear programming and PID controllers to manage budget pacing and bid optimization \cite{wu2015predicting, zhang2014real}. The field shifted significantly with the advent of Reinforcement Learning (RL), which framed the bidding process as a Markov Decision Process (MDP) \cite{cai2017real}. Recent SOTA approaches have leveraged Deep Q-Networks (DQN) and Proximal Policy Optimization (PPO) to handle the high-dimensional state space of auction features \cite{zou2019joint, feng2020deep}. However, these models often treat the market as a static environment or a slowly varying stochastic process, failing to account for the intentional, adversarial shifts in competitor behavior. 

Existing RL frameworks typically rely on explicit reward signals (e.g., Click-Through Rate or Return on Ad Spend) which are often delayed or sparse. This "reward-centric" view ignores the epistemic value of bidding actions—where a bid is not just an attempt to win, but a probe to uncover the valuation distribution of competitors \cite{ref_new_AdaptiveActiveI}. Furthermore, standard RL exploration strategies, such as $\epsilon$-greedy or entropy regularization, provide undirected stochasticity that is inefficient in the high-stakes, adversarial landscape of RTB where the "Winner's Curse" penalizes uninformed victories \cite{milosavljevic2020competition}.

\subsection{Active Inference in Sequential Decision Making}
Active Inference (AIF), rooted in the Free Energy Principle (FEP), offers a paradigm shift by casting both perception and action as an optimization problem: the minimization of Variational Free Energy (VFE) \cite{friston2010free}. Unlike RL, which maximizes an extrinsic reward, AIF agents minimize "surprise" relative to a generative model of the world \cite{friston2017active}. This inherently integrates perception, learning, and action. Recent advancements have demonstrated the utility of AIF in complex scenarios, such as autonomous reconnaissance missions where agents must "forage" for information to reduce uncertainty about hidden threats \cite{ref_new_ActiveInference}.

The implementation of AIF in high-dimensional spaces has been facilitated by Variational Autoencoders (VAEs) and modern variational message-passing schemes \cite{ref_new_RealizingSynthe, kingma2013auto}. Specifically, the use of Neural Ordinary Differential Equations (NODEs) allows for the continuous-time modeling of latent dynamics, which is crucial in RTB where auction frequencies occur at millisecond scales \cite{chen2018neural}. By treating the market as a continuous generative process, AIF agents can maintain a probabilistic belief over competitor intentions, allowing for proactive strategy adjustment rather than reactive response. This transition from external reward maximization to internal model-evidence maximization represents a critical evolution in agent design for non-stationary environments.

\subsection{Bridging Active Inference and Reinforcement Learning}
The theoretical distinction between AIF and RL is often debated, yet recent scholarship suggests that AIF subsumes RL by incorporating "the missing reward" into a broader framework of experience and epistemic foraging \cite{ref_new_TheMissingRewar}. While RL focuses on the pragmatic value (utility), AIF optimizes for both pragmatic and epistemic value through the minimization of Expected Free Energy (EFE) \cite{parr2019generalised}. In the context of RTB, the epistemic component of EFE compels the agent to bid in ways that resolve uncertainty about the "market price" and "competitor strength," effectively structuring the market through strategic information seeking.

Previous attempts to bridge these domains have looked at "RL as Inference" \cite{levine2018reinforcement} or Information Bottleneck theories \cite{tishby2000information}, but they often lack a principled mechanism for handling the Winner's Curse. The Winner's Curse occurs when the winner of an auction overestimates the item's value, a phenomenon exacerbated by information asymmetry. AIF provides a rigorous solution to this by explicitly modeling the hidden states of competitors as latent variables \cite{ref_new_AdaptiveActiveI}. By minimizing EFE, an agent naturally avoids bidding excessively in high-uncertainty regions unless the epistemic gain justifies the cost. This creates a "Strategic Market Structuring" effect, where the agent’s actions are motivated by a quest for a more predictable and controllable environment, leading directly into our proposed methodology of using AIF to navigate and mitigate the Winner's Curse.